\section{Conclusion}
Do we still need to be Agile? Yes, but Generative AI changes what effective
agility requires. Iteration, feedback, and adaptation remain valuable; however,
faster local implementation does not remove the need for shared architectural
intent, explicit interfaces, dependency boundaries, and reviewable acceptance
conditions. Across the three research questions, this study shows that
self-reported AI use is widespread and task-oriented, repository activity
concentrates near final integration, and observable upfront architectural
artifacts do not translate into a simple monotonic pattern of rework or
project outcomes. The late activity peak is not a single phenomenon: the
internal contrasts distinguish cases consistent with compensatory late
recovery, cases of late concentration without evaluator-score gain, and cases
of more distributed progress.

The study's contribution is to make this tension empirically inspectable through
complementary evidence: longitudinal perception and adoption reports, temporal
repository measures, transcript-based coordination evidence, structural T1
artifacts, path-based rework measures, and independent evaluator scorecards.
The robustness checks do not remove the observational limits of the study, but
they make the main validity threats inspectable through non-overlapping phase
bins, internal timing/outcome contrasts, operational-regularity proxies,
influence diagnostics, complexity profiles, and explicit transcript-coverage
boundaries. Together, these results suggest that GenAI-assisted development
shifts a larger share of engineering difficulty away from producing local code
and toward coordinating system-level intent. The practical implication is not a
return to immutable Big Design Up Front, but a stronger planning layer that
remains iterative, explicit, and shared by people and AI tools.

Spec-Driven Development is a plausible way to instantiate that planning layer:
requirements become specifications, specifications become technical plans and
tasks, and implementation is repeatedly checked against those artifacts. This
study does not evaluate SDD as an intervention, so its value here is a grounded
and testable proposition rather than an established causal conclusion. The
case for SDD is that it directly targets the coordination gap made visible by
the combined timing, planning, rework, complexity, and influence diagnostics.
Future work should compare ordinary practice with an iterative SDD condition
while measuring specification quality, consented AI-tool use, integration
timing, operational regularity, technical-complexity strata, coordination
evidence, rework provenance, and independently scored outcomes.

\section{Data Availability}
An anonymized replication package is available online~\cite{prism_replication_package}.
It contains the deterministic lexical-coding scripts, Git extraction queries,
and anonymized data needed to reproduce the metric derivations, with participant
privacy handled as described in Section~\ref{sec:data-sources}. Because the
evidence spans multiple sources and levels of analysis, we also provide an
interactive data explorer~\cite{prism_explorer} for filtering, navigating, and
inspecting the measures without manipulating the raw files directly.
